% Abhakseite „Das kann ich“ – Zone nur im Gesamt (\mitzone)
%% TODO Ich-kann-Sätze: Platzhalter wie in den Aufgabendateien
\begin{abhakseite}
\ifdefined\mitzone
\abhakgruppe{Kennst du schon}
\abhak{1}{Bruch in Dezimalzahl umwandeln (Division), abbrechende und periodische Dezimalzahlen erkennen und mit Periodenstrich schreiben, Dezimalzahl in Bruch}
\abhak{2}{Zahlengerade mit negativen Zahlen, Brüche und Dezimalzahlen eintragen, Zahlen vergleichen und ordnen}
\abhak{3}{Quadratwurzel als Umkehrung des Quadrierens, Quadratzahlen bis zwanzig hoch zwei, Wurzel mit dem Taschenrechner und Näherungswert, Wurzel zwischen zwei Nachbar-Quadratzahlen abschätzen, keine Wurzel aus negativer Zahl}
\abhak{4}{Taschenrechner}
\abhak{5}{Terme zusammenfassen}
\abhak{6}{Potenz als Malkette mit Basis und Exponent, Potenz ausrechnen (auch negative Basis), a hoch null gleich eins und negative Hochzahl als Bruch}
\abhak{7}{Brüche als Exponenten lesen und mit Brüchen rechnen}
\abhak{8}{Potenz als Malkette mit Basis und Exponent, Potenz ausrechnen (auch negative Basis), a hoch null gleich eins und negative Hochzahl als Bruch – Fehler finden}
\abhak{9}{Potenz als Malkette mit Basis und Exponent, Potenz ausrechnen (auch negative Basis), a hoch null gleich eins und negative Hochzahl als Bruch – selbst rechnen}
\fi
\abhakgruppe{Einheit 1 · Irrationale Zahlen und Zahlbereiche}
\abhak{10}{Zahlbereiche und irrationale Zahlen}
\abhak{11}{Dezimalzahl als Bruch schreiben (rational)}
\abhak{12}{Teilmengenkette ℕ ⊂ ℤ ⊂ ℚ ⊂ ℝ erklären (jede natürliche Zahl ist ganz, jede ganze rational, jede rationale reell)}
\abhak{13}{Zahlbereiche und irrationale Zahlen – Fehler finden, Begründen, Darstellungswechsel, Anwendung}
\abhakgruppe{Einheit 2 · Potenzgesetze}
\abhak{14}{Potenzgesetze}
\abhak{15}{Potenzen mit negativem Exponenten in Brüche und zurück (Blatt 0 aus potenzen-wurzeln.md; LISUM-PH)}
\abhak{16}{Potenzgesetze – Fehler finden, Begründen, Anwendung}
\abhakgruppe{Einheit 3 · Wurzelgesetze und rationale Exponenten}
\abhak{17}{Wurzelgesetze und rationale Exponenten}
\abhak{18}{Wurzelgesetze und rationale Exponenten – Fehler finden, Begründen, Darstellungswechsel, Anwendung}
\end{abhakseite}
